\chapter{Verification and certificates}\label{ch:verification}

\section{The ladder}

A program is certified by climbing six rungs. Each rung checks something the previous one cannot.

\begin{table}[h]
\centering\small
\begin{tabularx}{\textwidth}{@{}clX@{}}
\toprule
rung & establishes & how \\
\midrule
1 & sorts, shapes, feedback & \texttt{Program.check/1} \\
-- & exact rewriting & only identities valid over all of IEEE-754 \\
2 & admission & the allocation accepted by the extracted checker on \emph{every} instruction set; int8 contractions free of wrap-around at the maximum extents (\texttt{admissible/3}); SPIR-V limits \\
3 & adjoint identity & $\langle Wx, v\rangle = \langle x, W^{\mathsf T}v\rangle$ between the compiled kernel and the exact transpose: catches layout and transposition errors \\
4 & differential oracle & each CPU substrate equal to the exact oracle, bit for bit \\
5 & parity and envelope & equal FNV-1a digests across substrates; every output inside the envelope \\
6 & proof-carrying code & an Ed25519 certificate with a quorum of co-signatures \\
\bottomrule
\end{tabularx}
\caption{The verification ladder.}
\end{table}

\section{The certificate}

The payload holds the SHA-256 of the program, of each machine-code blob and of each SPIR-V module; the
digest of the Lean sources of the checkers; the evidence of each rung; and the \emph{counted work}
(floating-point operations, bytes, instructions per instruction set). It holds nothing that depends on
the host: no timings and no dispatch decision. Independent nodes that redo the ladder therefore produce
payloads identical byte for byte and can co-sign them (\texttt{Certificate.cosign/3});
\texttt{verify/3} requires a quorum of valid signatures from distinct trusted keys. A bundle packages
artefacts and certificate, and an edge node checks signatures and hashes in about two milliseconds and
runs the code without redoing the ladder.

\section{What is proved in Lean}

The Lean~4 development uses the core library only, with no \texttt{axiom} and no \texttt{sorry}. It
proves, among its 82 theorems:

\begin{itemize}
\item in $\mathbb{Z}/2^{32}\mathbb{Z}$: associativity and commutativity, that wrapping is a homomorphism,
  and that \emph{every} reduction tree with a wrapping adder computes the exact value when
  $K\cdot\max|A|\cdot\max|B| < 2^{31}$, for any permutation of the leaves (\texttt{integer\_parity});
  monotonicity of admissibility;
\item for memory banks: coprime strides are injective on the lanes (\texttt{bank\_injective}), and the
  closed-form padding for $2^m$ banks is sufficient and minimal;
\item for recurrences: the affine monoid is associative with neutral element $(1,0)$, composition is
  sequential application, and segmented lifting is associative for any monoid;
\item the soundness of the allocation checker (\texttt{checkAlloc\_sound});
\item the exact envelope decision and its monotonicity in the extent; Higham's bound on products of
  $(1+\delta_i)$;
\item the depth of Estrin's evaluation scheme against Horner's, computed over the tree.
\end{itemize}

\section{Extraction}

The extractor (\texttt{lake exe vapor-extract}) reads the \emph{elaborated} definitions (the same terms
the theorems speak about) and prints Elixir, failing on any construct outside its fragment. The
extracted functions are \texttt{check\_alloc}, \texttt{admissible}, \texttt{wrap\_s32},
\texttt{pad\_stride}, \texttt{bank\_of}, \texttt{affine\_op}, \texttt{seg\_affine} and
\texttt{within\_envelope}. Lean computes more than 480 conformance vectors that the extracted code must
reproduce, and the generated module carries the digest of the Lean sources, which a test checks without
needing the Lean toolchain. Extraction happens at build time: production never runs Lean.

\section{What is tested and what is trusted}

Tested but not proved: the encoders (against binutils and \texttt{spirv-val}, and by execution under
QEMU and lavapipe), the RVV interpreter against QEMU at VLEN 128, 256 and 512, bit-for-bit equality
across substrates, and fault containment. Trusted: the BEAM and the Linux kernel, the extractor's
printer and its six-line prelude, the IEEE-754 standard model, and the Vulkan driver, contained in its
own process.
